\section{Does It Generalize?}
\label{sec:generalization}
% DECISIONS.md s29 (pre-registration) and s30 (results). Phase A: three independent pre-registered
% families testing the model axis (logreg, MLP), the dataset axis (reprocessed CICIDS2017,
% UNSW-NB15), and the retention axis (sliding-window vs. accumulate). Full tables and per-cell
% statistics: docs/phaseA/{models,datasets,retention}-results.md. This section reports the
% verdicts and the numbers needed to read them, not the full tables.
Every result so far is on one real dataset (plus the synthetic one), two tree-ensemble models,
and one retention policy (\code{accumulate}). We pre-registered three generalization tests
before running them (DECISIONS.md \S29) — two more datasets, two weaker models with a different
inductive bias, and a sliding-window retention policy — and report what each found, including
where it did not replicate.

\begin{table*}[t]
  \centering
  \caption{Does each finding generalize? Columns: the original CICIDS2017 result (this paper,
  \S\ref{sec:degeneracy}--\S\ref{sec:attack}), the two new datasets (RF/XGBoost), and the two new
  models (on CICIDS2017). Key number in parentheses; full statistics in
  \texttt{docs/phaseA/\{models,datasets\}-results.md}.}
  % Source: results/phaseA/datasets/{cicids-fixed,unsw-nb15}/{nn_degeneracy_by_grid.csv,stats.csv},
  % results/phaseA/models/stats.csv, results/phaseA/retention/retention_stats.csv.
  \label{tab:generalization}
  \small
  \setlength{\tabcolsep}{2pt}
  \begin{tabular}{p{0.115\linewidth}p{0.175\linewidth}p{0.175\linewidth}p{0.175\linewidth}p{0.155\linewidth}p{0.125\linewidth}}
    \toprule
    Finding & CICIDS2017 (paper) & Reprocessed CICIDS2017 & UNSW-NB15 & Logreg & MLP \\
    \midrule
    Degeneracy & holds (benign NN 0.0065) & holds (0.0107) & holds (0.0037) & n/a & n/a \\
    Copy-fidelity FPR damage & holds for RF (+0.880 at $r{=}0.5$); XGBoost bimodal (+0.563 mean,
    3 of 5 seeds climb) & holds (RF +0.989, XGB +0.986) & holds (RF +0.254, XGB +0.211) & no
    (shows as TPR loss, $-0.10$ to $-0.42$) & partly ($r{=}0.5$ only, +0.409) \\
    Label-conflict isolation & holds (null arms flat) & partly (\code{a1} 3.3$\times$ the largest
    null arm) & doesn't (\code{a1truth} costs \emph{more} than \code{a1}) & partly (only
    \code{a1} lowers TPR) & doesn't \\
    Damage persists under a sliding window & holds (same final share) & not
    tested & not tested & not tested & not tested \\
    \bottomrule
  \end{tabular}
\end{table*}

\subsection{Degeneracy generalizes}
% DECISIONS.md s30.A/s30.B source numbers; results/phaseA/datasets/*/nn_degeneracy_by_grid.csv.
Both new datasets are as degenerate as CICIDS2017 under the same guard-(c) metric
(Table~\ref{tab:degeneracy-new}): every family with at least 500 \code{trusted_eval} rows sits at
least $10\times$ below the 0.25 leakage threshold, benign included. Reprocessing the flows (the
labeling and extraction fixes \citet{liu2022error} built) does not remove it — reprocessed
CICIDS2017's benign median is 0.0107, against the original release's 0.0065, both more than
$20\times$ below 0.25. The degeneracy appears under both CICFlowMeter features (both CICIDS2017
releases) and Argus/Bro features (UNSW-NB15), so it is a property of flow-summary features of
this kind of lab traffic, not of one tool or one release's extraction bugs.

\begin{table}[t]
  \centering
  \caption{Median NN RMS distance, \texorpdfstring{\code{trusted_eval}}{trusted\_eval} to the
  honeypot-pool-equivalent partition, grid 0.005, on the two new datasets. Compare
  Table~\ref{tab:degeneracy}'s CICIDS2017 row (BENIGN 0.0065).}
  % Source: results/phaseA/datasets/cicids-fixed/nn_degeneracy_by_grid.csv,
  % results/phaseA/datasets/unsw-nb15/nn_degeneracy_by_grid.csv.
  \label{tab:degeneracy-new}
  \small
  \begin{tabular}{lrr}
    \toprule
    Family & Reprocessed CICIDS2017 & UNSW-NB15 \\
    \midrule
    BENIGN      & 0.0107 & 0.0037 \\
    All attack  & 0.0079 & 0.0100 \\
    \bottomrule
  \end{tabular}
\end{table}

\subsection{Copy-fidelity FPR damage generalizes}
% DECISIONS.md s30.A/s30.B; results/phaseA/datasets/*/stats.csv (B.1, jitter 0, both ratios),
% results/phaseA/models/stats.csv (A.1, mlp, fixed dFPR).
A1's fixed-threshold FPR damage at copy-level fidelity clears Holm correction for RF and XGBoost
on both new datasets, at both poison ratios tested: reprocessed CICIDS2017 (RF +0.985/+0.989,
XGBoost +0.986/+0.986 at ratios 0.2/0.5) and UNSW-NB15 (RF +0.092/+0.254, XGBoost +0.016/+0.211).
The effect is not specific to CICIDS2017 or to tree ensembles. For the MLP it clears Holm at
ratio 0.5 only (+0.409, $p_t=0.00018$); ratio 0.2's effect (+0.0017) does not reach significance.
Logistic regression's A1 damage does not show up as fixed-threshold FPR at all — neither ratio
clears Holm under the paired $t$-test — but shows up instead as large, significant recalibrated-
threshold TPR loss at every cell tested (−0.10 to −0.42). The two new models relocate where the
damage surfaces; neither escapes it.

\subsection{Label-conflict isolation does not generalize}
% DECISIONS.md s30.A/s30.B B.2/A.2 results.
The mechanism-isolation test (\S\ref{sec:attack}: three matched null controls, each removing one
candidate cause) isolates the label conflict as the sole mechanism only on CICIDS2017 with RF —
every null arm's confidence interval includes 0 there. On every new dataset and model, at least
one null arm's 95\% CI excludes 0, so the isolation bar fails. The null-arm condition assumes a
detector near its ceiling, where extra correctly-labelled data has nowhere to push TPR; RF on
CICIDS2017 is at TPR $\approx 0.83$, but logistic regression (TPR $\approx 0.43$) and the MLP
(TPR $\approx 0.39$) are far from saturated, and so is RF on UNSW-NB15 (TPR $\approx 0.63$). Any
extra training data moves TPR there, label conflict or not, for reasons unrelated to the
mechanism the test is built to isolate.

Where the test cannot confirm isolation, the sign of the effect still carries some information.
On reprocessed CICIDS2017, \code{a1}'s effect (−0.051 $\Delta$TPR) is $3.3\times$ the largest
null arm's (\code{a1truth}, −0.015) — the label-conflict arm still does the most damage, even
though the nulls are not flat. On logistic regression, \code{a1} is the \emph{only} arm that
lowers TPR (−0.154); every correctly-labelled or uninformative arm \emph{raises} it
(\code{a1truth} +0.230, \code{s0j} +0.039, \code{junk} +0.144) — a qualitative pattern consistent
with the mechanism, reported as an observation, not a confirmed isolation. On UNSW-NB15 the
pattern reverses outright: \code{a1truth}, the correctly-labelled control, costs \emph{more} TPR
(−0.031) than \code{a1} itself (−0.021). The label conflict is not what the damage traces to
there.

\subsection{Damage does not need accumulation}
% DECISIONS.md s30.C/s30.C2; results/phaseA/retention/retention_stats.csv.
Everything above uses the loop's default retention, \code{accumulate} (every honeypot batch kept
forever). Tested against a sliding window instead, at a fixed final poison share, A1's damage is
unchanged: all 6 (retention $\times$ ratio) cells clear Holm, within 0.01 $\Delta$FPR of each
other at each ratio. Accumulation is not load-bearing for the damage this paper reports.
Section~\ref{sec:defenses} tests the harder case — a constant attacker budget, where a window
\emph{can} shrink the retained share — and what it implies for defending against this.

\subsection{Caveats}
UNSW-NB15 runs on 9 of the 24 canonical features (Argus/Bro flow fields, not CICFlowMeter's), and
its near-duplicate grid choice is not insensitive the way CICIDS2017's is — a coarser grid changes
RF's eval TPR from 0.74 to 0.52. Both new datasets run at $\tau=0$ (the s19 rule's own output for
each, not re-tuned to match CICIDS2017's $\tau=0.1$), so their at-distance ($j=0.7$) fixed-
threshold numbers are not directly comparable to the main paper's. Reprocessed CICIDS2017's
\code{trusted_eval} has no DDoS rows (a whole-burst partition artifact under real timestamps, not
a code defect). Logistic regression and the MLP run with the default hyperparameters in
\code{src/dloop/models/logreg.py} and \code{src/dloop/models/mlp.py}, not tuned to either
dataset.

\subsection{Pre-registration}
% DECISIONS.md s30, "Pre-registration problems found".
These generalization tests were pre-registered before any run (DECISIONS.md \S29). Reading the
results against that pre-registration surfaced four gaps in it, not in how the tests were run:
it does not say which significance test decides the damage condition (resolved as the paired
$t$-test, this paper's existing primary-test convention; the permutation reading alone would move
logistic regression's damage verdict from \`\`does not replicate'' to \`\`mixed,'' without changing the
label-conflict result either way); its null-arm bar assumes a saturated detector, which fails for
the weaker models and the new datasets for reasons unrelated to the mechanism it is meant to
isolate; its dataset-generalization bar combines two conditions evaluated at different
granularities (per-model and per-dataset), leaving XGBoost's verdict under-defined on both new
datasets; and its retention test used a budget mode under which the retention policy cannot
change the final poison share by construction, which a second, separately pre-registered family
(\S\ref{sec:defenses}) was run to correct.
